\section{Lab 3 - Luenberger Observer}

\subsection{Background Information}

In lab 3 the encoders are replaced by an IMU, consisting of a gyroscope that measures angular velocity in \(\un{rad/s}\) and an accelerometer that measures proper acceleration in \(\un{m/s^{2}}\).
Proper acceleration is measured relative to free fall, so a stationary accelerometer reads \(9.81\un{m/s^2}\) upwards.
The direction of gravity therefore gives the pitch and elevation angles \eqref{eq:4_1_Accelerometer}, see figure \ref{fig:Simplified_Helicopter_Forces_Angles}.
This only holds when the helicopter is not accelerating.

\begin{subequations}
    \begin{align}
        p & =\arctan\left( \dfrac{a_{y}}{a_{z}} \right) \label{eq:4_1_Pitch_From_Accelerometer} \\
        e & =\arctan\left( \dfrac{a_{x}}{\sqrt{ a_{y}^{2}+a_{z}^{2} }} \right) \label{eq:4_1_Elecvation_From_Accelerometer}
    \end{align} \label{eq:4_1_Accelerometer}
\end{subequations}

These angles are used to transform the gyroscope rates to \(\dot{p}\), \(\dot{e}\) and \(\dot{\lambda}\) using the equations given in the assignment \cite{helicopter_lab_assignment_2026}.

The IMU measurements are noisy, so instead of feeding them directly to the LQR controller, the states are estimated with a Luenberger observer \cite{chen_linear_system} based on the linearized model \eqref{eq:2_1_3_Lineariazed_Equations}, with
\(\mathbf{x}=\begin{bmatrix} p & \dot{p} & e & \dot{e} & \dot{\lambda} \end{bmatrix}^{\top}\),
\(\mathbf{u}=\begin{bmatrix} \tilde{V}_{s} & V_{d} \end{bmatrix}^{\top}\) and
\(\mathbf{y}=\begin{bmatrix} e & \dot{\lambda} \end{bmatrix}^{\top}\).
% TODO: Sjekk at dette var oppsettet under flyvningene (det er det i heli_q8_Part3.slx):
The gyroscope rates are low-pass filtered with a cut-off frequency of \(50\un{rad/s}\), and \(e\) is found by integrating the filtered elevation rate.
With this \(\mathbf{C}\), \verb|obsv(A,C)| has rank 5, so the system is observable without measuring every state.

The observer and the resulting estimation error \(\tilde{\mathbf{x}}=\mathbf{x}-\hat{\mathbf{x}}\) are
\begin{subequations}
    \begin{align}
        \dot{\hat{\mathbf{x}}} & =\mathbf{A}\hat{\mathbf{x}}+\mathbf{B}\mathbf{u}+\mathbf{L}(\mathbf{y}-\mathbf{C}\hat{\mathbf{x}}) \label{eq:4_1_State_Estimator} \\
        \dot{\tilde{\mathbf{x}}} & =(\mathbf{A}-\mathbf{L}\mathbf{C})\tilde{\mathbf{x}} \label{eq:4_1_Estimation_Error}
    \end{align}
\end{subequations}
so the error converges to zero when the eigenvalues of \(\mathbf{A}-\mathbf{L}\mathbf{C}\) lie in the left half-plane, and the pole placement decides how fast.
\(\mathbf{L}\) is found with \verb|L = place(A', C', P).'|.

\subsection{Hypothesis}

\begin{enumerate}[label=\textbf{H\arabic*}, leftmargin=*]
    \item \label{hyp:3_bias} At rest, the IMU measurements have small constant biases compared with the encoders. The observer assumes unbiased measurements, so these biases are not removed and will give offsets in the estimates.
    \item \label{hyp:3_accel} During flight, vibrations from the rotors and accelerations of the helicopter corrupt the accelerometer angles, so the measurement noise will be much larger than at rest.
    \item \label{hyp:3_noobs} Feeding the IMU measurements directly to the LQR controller passes all the measurement noise and errors to the controller, which will give worse control than with an observer.
    \item \label{hyp:3_fast} Observer poles far into the left half-plane give a large \(\mathbf{L}\). The estimates converge quickly but follow the measurements closely, so noise and biases pass through to the estimates and the controller.
    \item \label{hyp:3_slow} Observer poles close to the origin give a small \(\mathbf{L}\), so the estimates rely mostly on the model and lag behind the real states.
    By the separation principle, the closed-loop poles are the controller poles together with the observer poles \cite{chen_linear_system}, but because of model errors a slow observer in the loop will add lag that can cause oscillations or instability.
    \item \label{hyp:3_good} Observer poles moderately faster than the controller poles give estimates good enough to fly the helicopter with the joystick, but not as well as with the encoders.
\end{enumerate}

\subsection{Test Plans}

The encoder measurements are used as the true states in all tests.
In Tests 3--6 the estimates from the observer are used for feedback to the LQR controller from the previous lab.

\begin{table}[htbp]
    \centering
    \caption{Revised test plan for the Luenberger observer.}
    \label{tab:4_3_Test_Plans}
    \begin{tabularx}{\textwidth}{c X c X}
        \toprule
        \textbf{Test} & \textbf{Setup} & \textbf{Hyp.} & \textbf{Expected behavior} \\
        \midrule
        1 & IMU compared with encoders at rest and in flight & \ref{hyp:3_bias}, \ref{hyp:3_accel} & Small constant biases at rest, much more noise in flight \\
        2 & LQR fed directly by the IMU, no observer & \ref{hyp:3_noobs} & Poorer tracking and more oscillations \\
        3 & Observer poles \(-6,\dots,-10\), push during hover & \ref{hyp:3_good} & Stable hover, push is rejected \\
        4 & Observer poles around \(-100\) & \ref{hyp:3_fast} & Biases and noise reach the controller \\
        5 & Observer poles closer to the origin & \ref{hyp:3_slow} & Lag, oscillations or instability \\
        6 & Observer poles \(-6,\dots,-10\), flight with joystick & \ref{hyp:3_good} & The helicopter can be flown \\
        \bottomrule
    \end{tabularx}
\end{table}

\subsection{Results}

In Test 1 the IMU was compared with the encoders with the helicopter at rest and during flight, see table \ref{tab:4_4_IMU_Error}.
At rest the gyroscope has small constant biases, and the elevation from the accelerometer is \(0.55\un{rad}\) too high.
This offset comes from a constant of \(0.477\un{rad}\) that was added to the elevation in the Simulink model. Without this constant the offset was \(0.074\un{rad}\).
Since the elevation is used in the transformation of the gyroscope rates, the travel rate from the IMU was \(12\)--\(17\,\%\) too large during flight.
In flight the noise increased by a factor of \(70\)--\(150\), and the elevation from the accelerometer varied between \(0\) and \(1.2\un{rad}\) while the true elevation was close to zero, see figure \ref{fig:4_4_ElevationInFlight}.

\begin{table}[htbp]
    \centering
    \caption{Difference between IMU and encoder measurements at rest and in flight with observer poles \(-6,\dots,-10\). The mean is the bias and the standard deviation is the noise. The elevation includes an offset of \(0.477\un{rad}\) from the Simulink model.}
    \label{tab:4_4_IMU_Error}
    \begin{tabular}{l rr rr}
        \toprule
        & \multicolumn{2}{c}{At rest} & \multicolumn{2}{c}{In flight} \\
        \cmidrule(lr){2-3}\cmidrule(lr){4-5}
        Signal & Mean & Std & Mean & Std \\
        \midrule
        \(p\) [rad]               & \(-0.027\) & 0.001 & \(+0.006\) & 0.068 \\
        \(e\) [rad]               & \(+0.551\) & 0.002 & \(+0.586\) & 0.249 \\
        \(\dot{p}\) [rad/s]       & \(+0.007\) & 0.001 & \(+0.027\) & 0.075 \\
        \(\dot{e}\) [rad/s]       & \(-0.012\) & 0.002 & \(-0.010\) & 0.083 \\
        \(\dot{\lambda}\) [rad/s] & \(-0.004\) & 0.001 & \(+0.003\) & 0.042 \\
        \bottomrule
    \end{tabular}
\end{table}

\begin{figure}[htbp]
    \centering
    \includegraphics[width=\textwidth]{Images/3_Observer_Results/Report_ElevationInFlight.pdf}
    \caption{Elevation from the accelerometer compared with the encoder during flight with observer poles \(-6,\dots,-10\).}
    \label{fig:4_4_ElevationInFlight}
\end{figure}

Figure \ref{fig:4_4_ObserverPoles_Compare} compares the flights without the observer and with two different observer pole placements.
Without the observer (Test 2) the elevation stayed at about \(-0.18\un{rad}\), while it settled close to zero in the other runs, and pitch oscillated slowly at about \(0.1\un{Hz}\).
With the observer poles at \(-6,\dots,-10\) (Test 3) the elevation settled close to zero and the travel drift was small (\(0.04\un{rad/s}\)), but pitch kept oscillating at about \(0.45\un{Hz}\).
The push at \(t \approx 42\un{s}\) lowered the elevation by about \(0.13\un{rad}\), and the resulting pitch oscillation was damped within about \(6\un{s}\).
With the observer poles around \(-100\) (Test 4), roughly \(150\)--\(300\) times faster than the LQR poles at \(-0.3\) to \(-0.6\), the elevation was similar, but the travel drifted at \(-0.16\un{rad/s}\). % TODO: sjekk LQR-polene med Q og R fra kjøringene
In this run the IMU travel rate had a mean error of \(-0.05\un{rad/s}\), compared with \(0.003\un{rad/s}\) in Test 3.

\begin{figure}[htbp]
    \centering
    \includegraphics[width=\textwidth]{Images/3_Observer_Results/Report_ObserverPoles_Compare.pdf}
    \caption{Encoder measurements during hover with the LQR controller fed directly by the IMU (no observer), and with the Luenberger observer in the loop for two different pole placements. Travel is shown relative to its initial value.}
    \label{fig:4_4_ObserverPoles_Compare}
\end{figure}

With observer poles closer to the origin (Test 5) the helicopter could not take off, since it immediately rolled over to one side. % TODO: hvilke poler?
In all runs, also without the observer, the encoders show a constant pitch of \(-0.09\) to \(-0.13\un{rad}\), so the helicopter has a tendency to tilt to one side.
In the previous labs this was compensated by \textbf{[TODO: what did you do, e.g.\ an offset in \(V_d\)?]}, but this was not sufficient with the observer in the loop.
With the final observer the helicopter took off within about \(3\un{s}\), and it could be flown with the joystick for \(90\un{s}\) (Test 6), which agrees with \ref{hyp:3_good}.

\subsection{Discussion}

\paragraph{Fast and slow observer.}
A fast observer follows the measurements closely, so the bias in the travel rate is passed on to the estimates.
Since \(\ddot{\lambda}=K_{3}p\), an error in \(\dot{\lambda}\) gives an error in \(\hat{p}\), so the controller keeps the helicopter at the wrong pitch.
This is a possible explanation for the larger travel drift with poles around \(-100\), and it agrees with \ref{hyp:3_fast}.
A slow observer gave so much lag that the helicopter could not take off, which agrees with \ref{hyp:3_slow}.
Without the observer the elevation never reached zero, and pitch oscillated slowly, which agrees with \ref{hyp:3_noobs}.

\paragraph{Separation principle.}
Even with observer poles \(10\)--\(30\) times faster than the LQR poles, the helicopter oscillated in pitch, and slower observer poles made it impossible to take off.
The separation principle therefore does not hold perfectly in practice, because of model errors, the low-pass filters and the measurement errors.
The poles around \(-100\) are also faster than the low-pass filters at \(50\un{rad/s}\), so here the filters, and not the observer, decide how fast the estimates can follow the real states.

\paragraph{Measurement errors.}
The accelerometer angles were heavily corrupted during flight, as expected in \ref{hyp:3_accel}, and the constant gyroscope biases are not removed by the observer (\ref{hyp:3_bias}).
The elevation offset from the Simulink model also made the transformed travel rate \(12\)--\(17\,\%\) too large.
This is a model error that the observer cannot correct, and the constant should have been removed.

\paragraph{Physical meaning of the poles.}
A real observer pole \(\lambda\) gives a time constant \(\tau = 1/|\lambda|\) for the estimation error, so poles at \(-8\) correspond to \(\tau \approx 0.13\un{s}\) and poles at \(-100\) to \(\tau = 0.01\un{s}\).
However, the poles say nothing directly about how much noise reaches the estimates, which is why they were hard to choose.
This motivates a Kalman filter, where the gain is chosen from the noise statistics instead \cite{brown_kalman}.

\paragraph{Evaluation method.}
The observer was evaluated through the flight behaviour measured by the encoders.
Logging the estimates \(\hat{\mathbf{x}}\) as well would make it possible to compare the estimation error directly with the encoders.

\subsection{Conclusion}

The Luenberger observer made it possible to fly the helicopter using only IMU measurements, with clearly better results than feeding the IMU directly to the controller.
The pole placement is a trade-off between noise and lag: slow poles prevented take-off, while very fast poles let measurement biases through and gave a large travel drift.
Constant biases and model errors, such as the elevation offset, are not removed by the observer.
Since the poles cannot be chosen based on the actual noise levels, a Kalman filter is a natural next step.
